\section{总结与延伸}

\begin{enumerate}
    \item \textbf{Plan-Execute-RePlan 是一个非常 general 且有效的 Workflow}，适用于各种复杂任务场景，已被广泛验证。
    \item \textbf{自顶向下分解 + 自底向上反馈}：Plan 是"刚中之脑"的想象，Execute 是真实交互，RePlan 弥合两者之间的 gap。
    \item \textbf{经济性考量}：强模型做规划、弱模型做执行，与 OpenAI 的官方建议一致。
    \item \textbf{Plan 的质量决定上限}：如果任务复杂度高，可以在 Planning 阶段使用更强的模型或引入 Generator-Critique 机制。
    \item \textbf{从 Agentic Workflow 到 Autonomous Agent}：目前的实践证明，预定义的 Workflow 在可控性和可靠性上优于完全自主的 Agent。
\end{enumerate}

\end{document}
